% Part: lambda-calculus
% Chapter: introduction
% Section: lambda-definability

\documentclass[../../../include/open-logic-section]{subfiles}

\begin{document}

\olfileid{lam}{int}{rep}
\olsection{\usetoken{S}{lambda definable} حسابي تابعې}

لامبډا حساب څنګه د محاسبې نمونه کېدے شي؟ په لومړي سر کښې د دې
پوښتنې معنا هم روښانه نۀ ده۔ د طبيعي عددونو پر سر د محاسبه کېدنې د
څېړنې دپاره بايد د هغو عددونو مناسب تمثيل ومومو۔ لاندې يو داسې
تمثيل دے چې په حيرانوونکې توګه ښه کار کوي۔

\begin{defn}
د هر طبيعي عدد~$n$ دپاره \emph{د چرچ عددنښه} $\num{n}$ هغه لامبډا
ترم $\lambd[x][\lambd[y][(x(x(x(\dots x(y)))))]]$ تعريف کړئ چې
پکښې په ټوليز ډول $n$ ځله $x$ راغلے وي۔
\end{defn}

ترمونه $\num{n}$ ``تکراروونکي'' دي: پر ننوت~$f$، $\num{n}$ هغه
تابع ورکوي چې $y$،~$f^n(y)$ ته وړي۔ پام وکړئ چې هره عددنښه په عادي
بڼه کښې ده۔ اوس ويلے شو چې يو لامبډا ترم پر طبيعي عددونو يوه تابع
``محاسبه کوي'' يعنې څۀ۔

\begin{defn}
فرض کړئ $f(x_0, \dots, x_{k-1})$ له~$\Nat$ څخه~$\Nat$ ته يوه
$k$-ځاييزه جزوي تابع ده۔ وايو چې يو $\lambd$-ترم~$F$ هغه وخت، او
يوازې هغه وخت،~$f$ \emph{!!{lambda define}s} چې د طبيعي عددونو د
هرې لړۍ $n_0$، \dots،~$n_{k-1}$ دپاره،
\[
F\, \num{n_0}\, \num{n_1} \dots \num{n_{k-1}} \red \num{f(n_0, n_1, \dots,
  n_{k-1})}
\]
وي، که $f(n_0, \dots, n_{k-1})$ تعريف شوې وي؛ او که نۀ وي، نو
$F\, \num{n_0}\, \num{n_1} \dots \num{n_{k-1}}$ هېڅ عادي بڼه
نۀ لري۔
\end{defn}
\textbf{د ژباړې د نسخې سمون (OLLAM-002):} سرچينې د تابع آرګومېنټونه
د k په وسيله شمېرلي، خو په همدې تعريف کښې ئې تابع n-ځاييزه بللې ده۔
دلته ځاييزه کچه له هماغو k آرګومېنټونو سره برابره شوې ده۔
\textbf{د ژباړې د نسخې سمون (OLLAM-003):} سرچينې د ناتعريفه حالت
په تطبيق کښې د ترم او لومړۍ عددنښې تر منځ کامه ايښې وه؛ دلته هغه
د پاس ښودل شوي تطبيق په شان نېغ پرله‌پسې ليکل شوي دي۔

\begin{thm}
\ollabel{thm:lambda-def}
يوه تابع~$f$ هغه وخت، او يوازې هغه وخت، جزوي محاسبه کېدونکې تابع
ده چې يو لامبډا ترم ئې !!{lambda defined} کړي۔
\end{thm}

\begin{explain}
دا قضيه يو څه حيرانوونکې ده۔ د محاسبې د نمونې په توګه لامبډا حساب
يو ډېر ساده حساب دے؛ يوازې لامبډا تجريد او تطبيق ئې عملونه دي!{} خو
له دغو لږو وسيلو څخه هره محاسبوي کړنلاره پلې کېدے شي۔
\end{explain}

\end{document}
